\documentclass[10pt, letterpaper]{article}

% Packages:
\usepackage[
top=2 cm, % seperation between body and page edge from the top
bottom=2 cm, % seperation between body and page edge from the bottom
left=2 cm, % seperation between body and page edge from the left
right=2 cm, % seperation between body and page edge from the right
footskip=1.0 cm, % seperation between body and footer
% showframe % for debugging 
]{geometry} % for adjusting page geometry
\usepackage{titlesec} % for customizing section titles
\usepackage{tabularx} % for making tables with fixed width columns
\usepackage{array} % tabularx requires this
\usepackage[dvipsnames]{xcolor} % for coloring text
\definecolor{primaryColor}{RGB}{0, 0, 0} % define primary color
\usepackage{enumitem} % for customizing lists
\usepackage{fontawesome5} % for using icons
\usepackage{amsmath} % for math
\usepackage[
pdftitle={Jyotirmoy Avro's Cybersecurity Resume},
pdfauthor={Jyotirmoy Avro},
pdfcreator={LaTeX with RenderCV},
colorlinks=true,
urlcolor=primaryColor
]{hyperref} % for links, metadata and bookmarks
\usepackage[pscoord]{eso-pic} % for floating text on the page
\usepackage{calc} % for calculating lengths
\usepackage{bookmark} % for bookmarks
\usepackage{lastpage} % for getting the total number of pages
\usepackage{changepage} % for one column entries (adjustwidth environment)
\usepackage{paracol} % for two and three column entries
\usepackage{ifthen} % for conditional statements
\usepackage{needspace} % for avoiding page brake right after the section title
\usepackage{iftex} % check if engine is pdflatex, xetex or luatex

% Ensure that generate pdf is machine readable/ATS parsable:
\ifPDFTeX
\input{glyphtounicode}
\pdfgentounicode=1
\usepackage[T1]{fontenc}
\usepackage[utf8]{inputenc}
\usepackage{lmodern}
\fi

\usepackage{charter}

% Some settings:
\raggedright
\AtBeginEnvironment{adjustwidth}{\partopsep0pt} % remove space before adjustwidth environment
\pagestyle{empty} % no header or footer
\setcounter{secnumdepth}{0} % no section numbering
\setlength{\parindent}{0pt} % no indentation
\setlength{\topskip}{0pt} % no top skip
\setlength{\columnsep}{0.15cm} % set column seperation
\pagenumbering{gobble} % no page numbering

\titleformat{\section}{\needspace{4\baselineskip}\bfseries\large}{}{0pt}{}[\vspace{1pt}\titlerule]

\titlespacing{\section}{
% left space:
-1pt
}{
% top space:
0.3 cm
}{
% bottom space:
0.2 cm
} % section title spacing

\renewcommand\labelitemi{$\vcenter{\hbox{\small$\bullet$}}$} % custom bullet points
\newenvironment{highlights}{
\begin{itemize}[
topsep=0.10 cm,
parsep=0.10 cm,
partopsep=0pt,
itemsep=0pt,
leftmargin=0 cm + 10pt
]
}{
\end{itemize}
} % new environment for highlights


\newenvironment{highlightsforbulletentries}{
\begin{itemize}[
topsep=0.10 cm,
parsep=0.10 cm,
partopsep=0pt,
itemsep=0pt,
leftmargin=10pt
]
}{
\end{itemize}
} % new environment for highlights for bullet entries

\newenvironment{onecolentry}{
\begin{adjustwidth}{
0 cm + 0.00001 cm
}{
0 cm + 0.00001 cm
}
}{
\end{adjustwidth}
} % new environment for one column entries

\newenvironment{twocolentry}[2][]{
\onecolentry
\def\secondColumn{#2}
\setcolumnwidth{\fill, 4.5 cm}
\begin{paracol}{2}
}{
\switchcolumn \raggedleft \secondColumn
\end{paracol}
\endonecolentry
} % new environment for two column entries

\newenvironment{threecolentry}[3][]{
\onecolentry
\def\thirdColumn{#3}
\setcolumnwidth{, \fill, 4.5 cm}
\begin{paracol}{3}
{\raggedright #2} \switchcolumn
}{
\switchcolumn \raggedleft \thirdColumn
\end{paracol}
\endonecolentry
} % new environment for three column entries

\newenvironment{header}{
\setlength{\topsep}{0pt}\par\kern\topsep\centering\linespread{1.5}
}{
\par\kern\topsep
} % new environment for the header

\newcommand{\placelastupdatedtext}{% \placetextbox{<horizontal pos>}{<vertical pos>}{<stuff>}
  \AddToShipoutPictureFG*{% Add <stuff> to current page foreground
\put(
\LenToUnit{\paperwidth-2 cm-0 cm+0.05cm},
\LenToUnit{\paperheight-1.0 cm}
){\vtop{{\null}\makebox[0pt][c]{
\small\color{gray}\textit{Last updated in May 2025}\hspace{\widthof{Last updated in May 2025}}
}}}%
  }%
}%

% save the original href command in a new command:
\let\hrefWithoutArrow\href


\begin{document}
\newcommand{\AND}{\unskip
	\cleaders\copy\ANDbox\hskip\wd\ANDbox
	\ignorespaces
}
\newsavebox\ANDbox
\sbox\ANDbox{$|$}

\begin{header}
	\fontsize{25 pt}{25 pt}\selectfont Jyotirmoy Avro
	        
	\normalsize
	\kern 5.0 pt%
	\mbox{\hrefWithoutArrow{mailto:jyotirmoy.avro@gmail.com}{jyotirmoy.avro@gmail.com}}%
	\kern 5.0 pt%
	\AND%
	\kern 5.0 pt%
	\mbox{\hrefWithoutArrow{tel:+8801735342669}{(+880) 1735 342669}}%
	\kern 5.0 pt%
	\AND%
	\kern 5.0 pt%
	\mbox{\hrefWithoutArrow{https://github.com/itstheavro}{github.com/itstheavro}}%
	\kern 5.0 pt%
	\AND%
	\kern 5.0 pt%
	\mbox{\hrefWithoutArrow{https://linkedin.com/in/jyotirmoy-avro}{linkedin.com/in/jyotirmoy-avro}}%
\end{header}

\section{Objective} 
\begin{onecolentry}
	Highly motivated Computer Science student with practical cybersecurity experience from competitive Capture The Flag (CTF) events, seeking a remote Junior Security Analyst position. Eager to apply strong problem-solving skills, Python automation expertise, and a foundational understanding of network and web security to contribute to robust security operations and learn from industry professionals.
\end{onecolentry}

\section{Technical Skills}
\begin{onecolentry}
	\textbf{Cybersecurity Domains:} Web Application Security, Reverse Engineering, OSINT (Open Source Intelligence), IT Automation, Cryptography, Incident Response (Foundational)
\end{onecolentry}

\vspace{0.1 cm}

\begin{onecolentry}
	\textbf{Programming Languages:} Python (Scripting, Automation), JavaScript, C++, Java, PHP
\end{onecolentry}

\vspace{0.1 cm}

\begin{onecolentry}
	\textbf{Tools and Technologies:} Git, GitHub, Cisco Packet Tracer, Linux (Basic), VS Code, Jupyter Notebook, MySQL, Docker (Basic)
\end{onecolentry}
\vspace{0.1 cm}

\begin{onecolentry}
	\textbf{Concepts:} Machine Learning (for Security Analytics/Anomaly Detection), Data Structures \& Algorithms, Problem Solving
\end{onecolentry}

\section{Cybersecurity Competition Experience}
\begin{onecolentry}
	\begin{highlights}
		\item \textbf{BCS ICT FEST 2025 CTF:} Secured 15th place out of 162 teams in an intense 24-hour national-level Capture The Flag competition organized by Bangladesh Computer Society.
		\item \textbf{CyberRaid CTF 2025:} Ranked 16th out of 225 teams in the online round of this national-level competition, partnered with IT Directorate, Bangladesh Army, MIST, and KerberOS.
		\item \textbf{iutCTF 2024:} Achieved 15th place among 87 teams, successfully solving challenges in \textbf{PWN, Reverse Engineering, Web Exploitation, OSINT,} and \textbf{Machine Learning} categories.
		\item \textbf{SMP CTF Challenge 2024:} Secured 14th place out of 102 teams, qualifying for the final onsite round, demonstrating strong performance across diverse cybersecurity domains.
		\item \textbf{Inter University Capture The Flag (CTF) - BUET CSE Fest 2023:} Team PSTU\_NCREEPTION selected for the on-site round.
		\item \textbf{National Cyber Drill 2022:} Participated in the national cyber drill organized by BGD e-GOV CIRT (N-CERT).
	\end{highlights}
\end{onecolentry}

\section{Projects}

\begin{twocolentry}
	{\href{https://github.com/ItsTHEAvro/Cholesky-Key-Generator}{Code} | \href{https://itstheavro.github.io/Cholesky-Key-Generator/}{Live}}
	\textbf{Cholesky Key Generator}
\end{twocolentry}

\vspace{0.10 cm}
\begin{onecolentry}
	\begin{highlights}
		\item Developed a web-based cryptographic key generator using the Cholesky decomposition method.
		\item Demonstrates foundational understanding of cryptographic principles and secure key generation.
		\item Tools Used: HTML, CSS, JavaScript
	\end{highlights}
\end{onecolentry}

\vspace{0.2 cm}

\begin{twocolentry}
	{\href{https://github.com/ItsTHEAvro/ICMP_Pinger_Lab}{Code}}
	\textbf{ICMP Pinger Lab with Command Line and GUI Support}
\end{twocolentry}

\vspace{0.10 cm}
\begin{onecolentry}
	\begin{highlights}
		\item Developed a Python-based ICMP ping application measuring round-trip times and providing network latency statistics.
		\item Showcases foundational understanding of network protocols and network diagnostic capabilities crucial for network security.
		\item Tools Used: Python
	\end{highlights}
\end{onecolentry}

\vspace{0.2 cm}

\begin{twocolentry}
	{
		\href{https://codeforces.com/blog/entry/110724}{Codeforces Blog}
	}
	\textbf{Friends Online Status for Codeforces (Userscript)}
\end{twocolentry}

\vspace{0.10 cm}
\begin{onecolentry}
	\begin{highlights}
		\item Developed a JavaScript userscript to monitor friends' online status on Codeforces.
		\item Showcases practical scripting skills for web-based tools.
	\end{highlights}
\end{onecolentry}

\section{Certifications}
\begin{twocolentry}
	{\href{https://www.coursera.org/account/accomplishments/specialization/W6VVV4LRX95X}{Credential}}
	\textbf{Google IT Automation with Python Specialization (Google, February 2022)}
\end{twocolentry}
\vspace{0.10 cm}
\begin{onecolentry}
	\begin{highlights}
		\item Developed proficiency in Python programming and automation, critical for security scripting and operational efficiency.
		\item Skills: Python, Problem Solving, Automation
	\end{highlights}
\end{onecolentry}
\vspace{0.2 cm}

\begin{twocolentry}
	{\href{https://www.credly.com/badges/a1b2c3d4/public_url (replace with actual link if available)}{Credential}}
	\textbf{Introduction to Packet Tracer (Cisco, February 2024)} % Update month/year if different
\end{twocolentry}
\vspace{0.10 cm}
\begin{onecolentry}
	\begin{highlights}
		\item Gained foundational knowledge in network simulation and IoT, essential for understanding network security architectures.
		\item Skills: Cisco Packet Tracer, Network Simulation, IoT Simulation
	\end{highlights}
\end{onecolentry}
\vspace{0.2 cm}

\begin{twocolentry}
	{\href{https://www.coursera.org/account/accomplishments/specialization/BA8N9Y8JDSSN}{Credential}}
	\textbf{Machine Learning Specialization (DeepLearning.AI, April 2024)}
\end{twocolentry}
\vspace{0.10 cm}
\begin{onecolentry}
	\begin{highlights}
		\item Completed a comprehensive specialization in machine learning, applicable to security analytics, anomaly detection, and threat intelligence.
		\item Skills: Machine Learning, Python, Supervised/Unsupervised Learning, Deep Learning, Analytical Skills
	\end{highlights}
\end{onecolentry}
\vspace{0.2 cm}

\begin{twocolentry}
	{\href{https://freecodecamp.org/certification/itstheavro/javascript-algorithms-and-data-structures}{Credential}}
	\textbf{JavaScript Algorithms and Data Structures (freeCodeCamp, April 2023)}
\end{twocolentry}
\vspace{0.10 cm}
\begin{onecolentry}
	\begin{highlights}
		\item Mastered fundamental algorithms and data structures in JavaScript, enhancing capability for web application security analysis.
		\item Skills: Algorithms, Data Structures, JavaScript
	\end{highlights}
\end{onecolentry}
\vspace{0.2 cm}

\section{Education}
\begin{onecolentry}
	\textbf{Patuakhali Science and Technology University}, \textit{BSc. in Computer Science and Engineering}
\end{onecolentry}
\vspace{0.10 cm}
\begin{onecolentry}
	\textit{Year: 2022 - Present}\\ 
	\textit{Credits Completed: 104.75/165}\\
	\textit{CGPA: 3.567/4.00 (Position in class: 5\textsuperscript{th} out of 66 students)}\\
\end{onecolentry}
\section{Relevant Coursework}
\begin{onecolentry}
	\textbf{Computer Science:} Programming Language, Data Structure and Algorithms, Software Engineering, Database Systems, Artificial Intelligence, Machine Learning, Discrete Mathematics\\
	\vspace{0.2 cm}
	\textbf{Networking \& Communications:} Network Routing and Switching, Basic Communication Engineering, Optical Fiber Communication, and related lab courses.\\
	\vspace{0.2 cm}
	\textbf{Related Technical Fields:} Digital Logic Design, Microprocessor and Interfacing, Basic Electrical Engineering, Electrical Technology.
\end{onecolentry}

\end{document}